\documentclass{article}

\usepackage[nonatbib, final]{neurips}
\usepackage[numbers]{natbib}

\makeatletter
\renewcommand{\@noticestring}{
  \centering
  Mid-Project Report
}
\makeatother

\input{extra_pkgs}
\usepackage{physics}
\usepackage{mathtools}
\DeclarePairedDelimiter\p{(}{)}
\DeclarePairedDelimiter\n{|}{|}
\DeclarePairedDelimiter\B{[}{]}

% ── Custom macros ─────────────────────────────────────────────────────────────
\newcommand{\ShiftDC}{\textsc{ShiftDC}}
\newcommand{\sss}{\textsc{sss}}
\newcommand{\ssu}{\textsc{ssu}}
\newcommand{\xvl}{x_{\text{vl}}}
\newcommand{\xtt}{x_{\text{tt}}}
\newcommand{\ml}{m^l}
\newcommand{\safel}{s^l}

\title{Modality-Conditioned Safety Representations\\for Vision-Language Models\\[4pt]
       \large Mid-Project Report}

\author{%
  % Authors
}

\begin{document}

\maketitle

% ─────────────────────────────────────────────────────────────────────────────
\begin{abstract}
Vision-Language Models (VLMs) exhibit a subtle safety failure: given a \emph{safe image}
paired with a \emph{safe text prompt}, the model sometimes produces an unsafe response---a
phenomenon we term \textsc{ssu} (Safe-image, Safe-text $\rightarrow$ Unsafe output).
We argue this failure calls for \textbf{modality-conditioned safety representations}: safety
boundaries that adapt based on how much each modality contributes to the input.
As a first step, we apply the \ShiftDC{} diagnostic~\cite{zou2025shiftdc} to LLaVA-1.5-7B
on HoliSafe-Bench~\cite{lee2025holisafe} (682 \textsc{sss} + 718 \textsc{ssu} samples),
computing a per-layer modality shift $\ml = \xvl^l - \xtt^l$ and projecting it onto a
PCA-derived safety direction $\safel$.
We find significant \textsc{sss}/\textsc{ssu} divergence in 25 of 32 layers ($p < 0.001$),
with the largest gap at layer 31 ($p = 5.5\times10^{-30}$).
A text-only sanity check reveals that \textsc{ssu} prompts are already semantically closer
to unsafe content before image integration, suggesting the failure is partly prompt-driven.
Building on these diagnostic findings, we are developing a modality-conditioned correction
mechanism guided by attention allocation signals.
\end{abstract}

% ─────────────────────────────────────────────────────────────────────────────
\section{Introduction}
% ─────────────────────────────────────────────────────────────────────────────

Vision-Language Models (VLMs) extend large language models with visual perception, enabling
reasoning over image-text pairs~\cite{liu2024llava15}.
Despite strong performance on downstream tasks, VLMs introduce safety vulnerabilities
qualitatively distinct from their text-only counterparts.
Prior work has shown that safety-relevant concepts in large language models are encoded as
linear directions in activation space~\cite{arditi2024refusal}, enabling effective safety classification through geometric methods.
However, VLMs face a more complex challenge: visual inputs induce systematic shifts in these
safety representations~\cite{zou2025shiftdc}, causing the model to perceive certain
multimodal inputs as \emph{safer} than equivalent text-only queries.

A striking manifestation of this distortion is the \textbf{\textsc{ssu} phenomenon},
identified in HoliSafe-Bench~\cite{lee2025holisafe}: given a \emph{safe image} and a
\emph{safe text prompt}, the model nonetheless produces an unsafe response.
Neither modality alone triggers the failure, and text-only safety mechanisms
(RLHF, content classifiers) are blind to it.

We argue that effective VLM safety detection requires \textbf{modality-conditioned safety
representations}---boundaries that adapt based on the dominant modality contributing to a
given input.
Our project proceeds in two phases:
\begin{enumerate}[noitemsep]
  \item \textbf{Diagnostic phase (complete):} Characterize \textsc{sss}/\textsc{ssu}
        activation differences layer-by-layer using the \ShiftDC{} framework.
  \item \textbf{Intervention phase (in progress):} Learn modality-conditioned safety
        subspaces guided by the diagnostic findings.
\end{enumerate}

% ─────────────────────────────────────────────────────────────────────────────
\section{Related Work}
% ─────────────────────────────────────────────────────────────────────────────

\paragraph{Linear representation of safety.}
Arditi et al.~\cite{arditi2024refusal} demonstrated that refusal behavior in LLMs is mediated
by a \emph{single} linear direction in activation space, and that ablating this direction
reliably suppresses refusal across diverse harmful inputs.
This geometric structure underpins our use of a single per-layer safety direction $\safel$.

\paragraph{Modality-induced safety distortion.}
\ShiftDC{}~\cite{zou2025shiftdc} shows that visual inputs cause consistent activation shifts
that distort safety perception: multimodal inputs appear safer than text-only equivalents in
hidden-state geometry even when semantic content is held fixed.
The authors propose a global projection-based correction; a key limitation is that this
correction is modality-agnostic.
Our work addresses this by conditioning safety boundaries on the dominant modality.

\paragraph{Attention imbalance in VLMs.}
Recent mechanistic studies have observed that text tokens attend to image tokens far less
than their proportion of the input would suggest, contributing to safety vulnerabilities.
This attention imbalance motivates using the attention allocation toward image tokens as a
lightweight modality conditioning signal in our Phase~2 classifier.

\paragraph{Benchmarks.}
HoliSafe-Bench~\cite{lee2025holisafe} provides per-modality safety labels across five
image-text combinations, making it uniquely suited for \textsc{sss}/\textsc{ssu} differential
analysis.
MM-SafetyBench~\cite{liu2024mmsafetybench} and
JailBreakV-28K~\cite{luo2024jailbreakv} target adversarial attacks.


% ─────────────────────────────────────────────────────────────────────────────
\section{Technical Approach}
% ─────────────────────────────────────────────────────────────────────────────

\subsection{Phase 1: Activation Extraction and Safety Direction}

For each sample in HoliSafe-Bench we extract hidden states at every transformer layer via
two forward passes through LLaVA-1.5-7B:
\begin{itemize}[noitemsep]
  \item \textbf{VL activations} $\xvl^l \in \mathbb{R}^d$: last-token hidden state from
        the full multimodal input (image + text prompt).
  \item \textbf{TT activations} $\xtt^l \in \mathbb{R}^d$: last-token hidden state from
        a text-only surrogate (generated image caption + text prompt).
\end{itemize}
The \textbf{modality shift} $\ml = \xvl^l - \xtt^l$ isolates the contribution of visual
integration.
Captions are generated by LLaVA-1.5-7B (batched, \texttt{max\_new\_tokens=100}).
All activations are stored via an on-disk \texttt{ActivationCache} for reproducibility.

\paragraph{PCA-based safety direction.}
We compute $\safel$ from two reference datasets: LLaVA-Instruct-80k (safe) and
MM-SafetyBench~\cite{liu2024mmsafetybench} (unsafe).
Rather than the simple mean difference (which confounds safety with semantic content of
the reference corpora), we use a PCA-based steering vector~\cite{zou2025shiftdc}:
\begin{enumerate}[noitemsep]
  \item Collect $H^+ \in \mathbb{R}^{N^+ \times d}$ (safe) and
        $H^- \in \mathbb{R}^{N^- \times d}$ (unsafe).
  \item Jointly mean-center: $\mu = \tfrac{1}{2}(\bar H^+ + \bar H^-)$,
        $M = [H^+ - \mu;\; H^- - \mu]$.
  \item $\safel$ = first right singular vector of $M$, sign-flipped to align with
        $\bar H^+ - \bar H^-$ (unsafe $\to$ safe).
\end{enumerate}

\paragraph{Shift projection.}
For each sample $i$ and layer $l$:
\begin{equation}
  \text{cosine}^l_i = \cos(\ml_i,\,\safel), \qquad
  \text{proj}^l_i   = \frac{\ml_i \cdot \safel}{\|\safel\|^2}.
\end{equation}
Positive values indicate the modality shift is directed toward the \emph{safe} side.
We compare \textsc{sss} vs.\ \textsc{ssu} groups per layer via Welch's $t$-test.

\subsection{Phase 2: Modality-Conditioned Safety Classifier (In Progress)}

Building on the diagnostic findings, we are developing a modality-conditioned safety
boundary.
We use attention allocation as a lightweight conditioning signal.
Let $\alpha_\text{img}^{(l)}$ be the fraction of attention mass on image tokens at layer $l$.
Define the modality score:
\begin{equation}
  s_\text{mod}(\mathbf{x}) = \frac{1}{L}\sum_{l=1}^{L}\alpha_\text{img}^{(l)},
\end{equation}
and interpolate between two learned projection matrices:
\begin{equation}
  \mathbf{P}(\mathbf{x}) =
    (1-s_\text{mod})\,\mathbf{P}_\text{txt} + s_\text{mod}\,\mathbf{P}_\text{img}.
\end{equation}
Safety is classified as
$f(\mathbf{x}) = \operatorname{sign}\!\bigl(\mathbf{P}(\mathbf{x})\,h^{l^*}\cdot\mathbf{w}+b\bigr)$,
where the critical layer $l^*$ is selected via the diagnostic results (Section~\ref{sec:results}).
We jointly learn $\mathbf{P}_\text{txt}$ and $\mathbf{P}_\text{img}$ with a contrastive
objective:
\begin{equation}
  \mathcal{L} = \mathcal{L}_\text{contrastive}
              + \lambda\,\|\mathbf{P}_\text{txt} - \mathbf{P}_\text{img}\|_F,
\end{equation}
where the regularizer encourages the two projections to differ meaningfully.

% ─────────────────────────────────────────────────────────────────────────────
\section{Progress \& Preliminary Results}
\label{sec:results}
% ─────────────────────────────────────────────────────────────────────────────

\paragraph{Setup.}
Model: LLaVA-1.5-7B~\cite{liu2024llava15} ($d = 4096$, 32 transformer layers).
Dataset: HoliSafe-Bench, filtered to \textbf{682 \textsc{sss}} + \textbf{718 \textsc{ssu}}
samples.
Reference datasets for safety direction: LLaVA-Instruct-80k (safe) and MM-SafetyBench
(unsafe).
Phase 1 is fully implemented as a reproducible pipeline (\texttt{run\_shiftdc.sh})
with modular GPU extraction and CPU analysis stages.

\subsection{Main \ShiftDC{} Results}

\begin{table}[t]
  \centering
  \caption{Per-layer \ShiftDC{} cosine similarity results (selected layers).
    Positive gap = \textsc{sss} $>$ \textsc{ssu}.
    $p$-values from Welch's $t$-test (two-sided).}
  \label{tab:shiftdc}
  \begin{tabular}{ccccc}
    \toprule
    Layer & \textsc{sss} cosine & \textsc{ssu} cosine & Gap & $p$-value \\
    \midrule
    1  & 0.075 & 0.072 & $-$0.002 & 0.188 (n.s.) \\
    4  & 0.452 & 0.456 & $-$0.004 & $6.8\times10^{-4}$ \\
    8  & 0.260 & 0.283 & $-$0.022 & $3.0\times10^{-20}$ \\
    12 & 0.041 & 0.068 & $-$0.027 & $3.5\times10^{-15}$ \\
    20 & 0.057 & 0.039 & $+$0.018 & $1.0\times10^{-9}$ \\
    \textbf{31} & \textbf{0.092} & \textbf{0.055} & $+$\textbf{0.037} &
      $\mathbf{5.5\times10^{-30}}$ \\
    \bottomrule
  \end{tabular}
\end{table}

Table~\ref{tab:shiftdc} shows representative results across 32 layers.
\textsc{sss} and \textsc{ssu} groups show statistically significant differences
($p < 0.001$) in \textbf{25 of 32 layers}.
Two regimes emerge:

\textbf{Early-to-mid layers (4--14):} \textsc{ssu} samples exhibit a \emph{larger} positive
modality shift toward the safe direction.
This may indicate that the model attempts to regulate its internal representation as the
unsafe potential of these samples is not yet fully represented.

\textbf{Late layers (20--31):} The pattern reverses---\textsc{sss} samples show a larger
safe-direction projection.
The maximum gap at \textbf{layer 31} (\textsc{sss} = 0.092, \textsc{ssu} = 0.055,
$p = 5.5\times10^{-30}$) indicates that \textsc{ssu} final representations sit closer to
the unsafe side of $\safel$, consistent with the unsafe output.
This identifies layer 31 as a strong candidate for $l^*$ in the Phase 2 classifier.

\subsection{Sanity Check: Text-Only Baseline}

This experiment tests whether the \textsc{sss}/\textsc{ssu} difference is caused by image
integration or by pre-existing semantic differences in the prompts.

\textbf{Finding:} In \textbf{28 of 32 layers}, \textsc{ssu} caption-only (TT) activations
project more strongly onto the \emph{unsafe} side of $\safel$ than \textsc{sss} TT activations.
For example, at Layer 22: \textsc{sss} projection = $-5.09$, \textsc{ssu} projection =
$-13.82$.

This indicates that \textsc{ssu} prompts are semantically closer to unsafe content
\emph{before} any image is processed.
The larger positive modality shifts seen in early layers (Section 4.1) likely reflect this
baseline difference rather than purely image-induced distortion.
This finding motivates the modality-conditioning approach: a safety boundary that ignores
the text-level baseline risks either over-triggering on benign multimodal inputs or
missing failures driven by prompt semantics.

\subsection{Experiment A: Effective Rank of Modality Shift Space}

The modality shift occupies a \textbf{low-dimensional subspace} relative to $d = 4096$.
At layer 2 with cumulative variance threshold $\tau = 0.9$, effective rank $\approx 88$--$96$
for both \textsc{sss} and \textsc{ssu}.
Crucially, the two groups show nearly identical rank profiles across all layers, indicating
the \textsc{sss}/\textsc{ssu} difference is in the \emph{direction} of the shift, not its
dimensionality.
The low rank ($\ll d$) confirms that safety-relevant distortions are structured and suggests
that a low-rank projection $\mathbf{P}$ in Phase 2 is well-motivated.

% ─────────────────────────────────────────────────────────────────────────────
\section{Challenges \& Risks}
% ─────────────────────────────────────────────────────────────────────────────

\paragraph{Text-level confound in \textsc{ssu} samples.}
The sanity check shows that \textsc{ssu} prompts are semantically closer to unsafe content
in text space.
Attributing the failure purely to image integration would be incorrect, and
disentangling the two effects requires controlled ablations not yet completed.

\paragraph{Layer selection sensitivity.}
Safety directions vary substantially across transformer layers.
The diagnostic analysis identifies layer 31 as the most discriminative, but selecting
$l^*$ robustly for the Phase 2 classifier requires careful cross-layer validation to
avoid overfitting.

\paragraph{Semantic confound in reference datasets.}
The safety direction $\safel$ is computed from LLaVA-Instruct vs.\ MM-SafetyBench, which differ
not only in safety-relevance but also in topic distribution and writing style.
The PCA method reduces this confound but does not eliminate it.
We plan to introduce a paired text-only dataset (\texttt{safeunsafeds}) with more controlled
distributions.

\paragraph{Generalization of linear interpolation.}
The attention-based modality conditioning assumes a smooth spectrum between text-dominant
and image-dominant regimes.
For \textsc{ssu} inputs---where both modalities are safe individually---the conditioning
signal may be ambiguous, potentially requiring a non-linear conditioning approach.

\paragraph{Compute constraints.}
Running two forward passes per sample across 32 layers on a 7B-parameter model is
memory-intensive.
Extending to larger models (e.g., LLaVA-1.5-13B) or secondary benchmarks
(MM-SafetyBench, JailBreakV-28K) requires additional GPU time.

% ─────────────────────────────────────────────────────────────────────────────
\section{Next Steps \& Timeline}
% ─────────────────────────────────────────────────────────────────────────────

\begin{table}[h]
  \centering
  \caption{Remaining project milestones.}
  \label{tab:timeline}
  \begin{tabular}{lp{9cm}}
    \toprule
    Dates & Milestone \\
    \midrule
    Mar 31 -- Apr 6
      & Complete follow-up experiments (B: safety decomposition, C: subspace overlap,
        D: category analysis); determine whether \textsc{ssu} distortion is
        concentrated in $\safel$ or spread across multiple dimensions. \\
    Apr 7 -- Apr 13
      & Integrate paired \texttt{safeunsafeds} reference dataset; recompute $\safel$ and
        compare to LLaVA-Instruct/MM-SafetyBench baseline.
        Begin Phase 2 implementation: extract attention allocation signals and
        implement interpolated projection classifier. \\
    Apr 14 -- Apr 20
      & Evaluate Phase 2 classifier on HoliSafe-Bench \textsc{ssu} vs.\ \textsc{sss};
        compare against the ShiftDC~\cite{zou2025shiftdc} baseline.
        Ablate layer choice, $\lambda$, and linear vs.\ nonlinear conditioning. \\
    Apr 21 -- Apr 27
      & Extend analysis to LLaVA-1.5-13B to test generality; produce PCA visualizations
        and error analysis for final report. \\
    Apr 28 -- May 4
      & Final report revisions and submission. \\
    \bottomrule
  \end{tabular}
\end{table}

The core diagnostic pipeline (Phase 1) is fully implemented and reproducible.
The primary remaining work is: (1) completing follow-up experiments B--D to characterize
the geometry of safety distortion; (2) addressing the semantic confound through a cleaner
reference dataset; (3) implementing and evaluating the Phase 2 modality-conditioned
classifier; and (4) validating findings on a second model.

% ─────────────────────────────────────────────────────────────────────────────
\bibliographystyle{plainnat}
\bibliography{bibliography}

\end{document}
